\documentclass[12pt]{article}

\usepackage[margin=1in]{geometry}
\usepackage{setspace}
\usepackage{amsmath,amssymb,amsthm}
\usepackage{graphicx}
\usepackage{booktabs}
\usepackage{array}
\usepackage{multirow}
\usepackage{threeparttable}
\usepackage{natbib}
\usepackage{hyperref}
\usepackage[table]{xcolor}
\usepackage{caption}
\usepackage{float}
\usepackage{placeins}
\usepackage{enumitem}
\usepackage{titlesec}
\usepackage{siunitx}

\doublespacing
\hypersetup{colorlinks=true,linkcolor=black,citecolor=black,urlcolor=blue}
\captionsetup{font=small,labelfont=bf,justification=raggedright,singlelinecheck=false}
\setlength{\parindent}{0.5in}
\setlength{\parskip}{0pt}

\titleformat{\section}[block]{\normalfont\normalsize\bfseries\centering}
  {\Roman{section}.}{0.5em}{\MakeUppercase}
\titleformat{\subsection}[block]{\normalfont\normalsize\itshape}
  {\Alph{subsection}.}{0.5em}{}
\titleformat{\subsubsection}[runin]{\normalfont\normalsize\itshape}
  {}{0em}{}[.---]
\titlespacing{\section}{0pt}{18pt}{6pt}
\titlespacing{\subsection}{0pt}{12pt}{4pt}

\definecolor{darkblue}{RGB}{10,30,80}

% ─────────────────────────────────────────────────────────────────────────────
\begin{document}

% ── Title Page ────────────────────────────────────────────────────────────────
\begin{titlepage}\begin{center}\vspace*{1.2cm}

{\Large\bfseries Fire Without Smoke:\\[8pt]
Geopolitical Risk and the New Inflation Decoupling}

\vspace{1.4cm}
{\normalsize Eleni Kalamara%
\footnote{King's College London and Nomura Securities International.
E-mail: \texttt{eleni.kalamara@kcl.ac.uk}.
Replication materials: \texttt{https://github.com/Elenikal/geoshock}.
Live dashboard: \texttt{https://geoshock.onrender.com}.
The author thanks Claude (Anthropic) for research assistance.
The views expressed are the author's own.
All errors are the author's own.}}

\vspace{0.9cm}{\normalsize March 2026}\vspace{1.4cm}

\begin{abstract}\noindent
Since 2023, the US has experienced sustained crisis-level geopolitical risk---peaking
at $z = 5.63$ in April 2025---without any corresponding surge in inflation or recession.
We explain why by introducing the Geopolitical Inflation Pressure Index (GIPI),
which measures whether the economic channels through which geopolitical shocks
transmit to prices are open or closed, and embedding it in a
Growth-at-Risk framework via a GPR$\,\times\,$GIPI interaction term.
The amplification hypothesis is confirmed robustly in-sample: GIPI
carries significant information across quantile-horizon
cells---amplifying left-tail risk at $\tau \leq 0.25$ and capturing
energy-sector windfalls at upper quantiles---with the largest
tail-risk gains at long horizons.
The 2023--25 GPR--GIPI decoupling ($\rho = 0.131$) reflects US shale
supply, IEA reserve coordination, and Fed credibility closing all five
transmission channels simultaneously.
A real-time pipeline---incorporating VLCC freight rates
(TD3C route) and a headline-override mechanism---confirms the thesis
live: as of mid-March 2026, elevated GPR coexists with subdued GIPI
and an inactive AIS trigger.
The unfolding closure of the Strait of Hormuz (late March~2026)
provides an out-of-sample real-time stress test, demonstrating
the framework's capacity to detect regime transitions as
transmission channels begin to reopen.
\end{abstract}

\end{center}\end{titlepage}

\newpage\setcounter{page}{1}

\noindent\textbf{JEL Classification:} E27, E32, F51, Q43, G17\\[4pt]
\noindent\textbf{Keywords:} geopolitical risk, growth-at-risk, quantile regression,
inflation transmission, local projections, GDELT, real-time monitoring

\bigskip

% ═══════════════════════════════════════════════════════════════════════════════
\section{Introduction}
% ═══════════════════════════════════════════════════════════════════════════════

Between March and November 2025, the Geopolitical Risk index of
\citet{caldara2022} averaged $z = 3.26$ standard deviations above its
historical mean---well into crisis territory, with an April peak of
$z = 5.63$ driven by the US--Israel--Iran military exchange.
Yet US industrial production grew 1.3\% YoY in December 2025,
CPI inflation stood at 2.65\%, and the six-month severe-contraction probability
implied by the standard Growth-at-Risk model was 6.4\%.\footnote{%
  Throughout the paper, ``recession probability'' refers to
  $\Pr(\text{IP YoY} < -2\%)$, a severe-contraction threshold broadly
  consistent with NBER recession episodes in our sample; IP growth alone
  does not determine NBER dating.}
Something fundamental had changed.

The conventional transmission story runs as follows.
A Middle East conflict disrupts oil supply, raising Brent crude, which
passes through to headline CPI within 6--9 months \citep{conflitti2019}.
Rising energy costs squeeze margins, forcing firms to cut output and
hiring \citep{bloom2009}.
Simultaneously, import prices surge as shippers reprice Red Sea routing
risk, food prices spike as grain supply chains re-route, and inflation
expectations de-anchor \citep{caldara2025}.
The result---as observed after the February 2022 Russian invasion
of Ukraine---is a joint supply shock across all five channels.
The invasion disrupted European energy supply, rerouted Black Sea
grain shipments, and broke global supply chains at a moment when
US inflation was already running at 7.5\% YoY on the back of
COVID-era fiscal stimulus and supply-chain dislocation
\citep{bernanke2023}; the additional shock across all five
channels drove GIPI from $-0.5$ in late 2021 to $+5.6$ by
June 2022, amplifying CPI to its highest level since 1981.\footnote{%
  US CPI peaked at 9.1\% YoY in June 2022.
  Pre-invasion (January 2022) CPI was already 7.5\%, driven by
  pandemic-era demand-supply imbalances; Ukraine added the energy
  and food-price leg that pushed it to the peak.}
The contrast with 2025 could not be starker: GPR reached similar
peak levels in both episodes, yet GIPI in 2025 remained near zero.

Why was 2025 different?
All five channels failed to engage.
Global energy prices fell 14--17\% YoY throughout 2025 even as
GPR reached record highs, because US crude production stood at
record levels of 13.5 million barrels per day \citep{eia2025},
IEA members coordinated strategic reserve releases within days of
each escalation announcement \citep{aiyar2023}, and
the Federal Funds rate at 3.72\% kept TIPS breakevens anchored
near 2.3\%.
GIPI, constructed as the first principal component of the five
channels, read $-1.01$ in December 2025---a full 6.6 GIPI units
below its 2022 peak.
The GPR--GIPI correlation over these 35 months was $\rho = 0.131$.

This paper identifies the conditions under which geopolitical risk
transmits to macro tail risk: conditioning on whether inflation
transmission channels are open or closed---as measured by GIPI---
substantially changes the in-sample tail of the IP growth distribution
relative to a standard GPR-based GaR model.
The decoupling (US shale, IEA reserves, Fed credibility) and the
Layer~0 real-time surveillance system are complementary contributions
that motivate and operationalise the framework.
Causal language (``driven by,'' ``amplified by'') describes the
reduced-form state-dependence the data reveal, not identified structural
parameters; the frontier extensions that would sharpen causal
identification are discussed in Section~\ref{sec:conclude}.

\medskip
This paper makes four contributions to the literatures on geopolitical risk,
growth-at-risk, and real-time monitoring.

\textit{First}, we introduce GIPI---the Geopolitical Inflation Pressure Index---as
the first principal component of five inflation-transmission-channel series.
Unlike the GPR index, GIPI measures whether the economic pipes through which
geopolitical shocks travel to prices are open or closed.

\textit{Second}, we augment the Growth-at-Risk framework with a GPR$\times$GIPI
interaction term, allowing the tail of the IP growth distribution to respond
differently to GPR depending on channel stress.
The amplification hypothesis is confirmed in-sample across 16 of 18
quantile-horizon cells shown in Table~\ref{tab:orth}, with the largest
gains at the long-horizon left tail.

\textit{Third}, we document the 2023--25 GPR--GIPI decoupling as a novel
macroeconomic stylised fact and provide a structural interpretation grounded
in US shale supply, IEA reserve coordination, and Fed credibility.
The decoupling explains why standard GPR-based models overpredicted tail risk
throughout 2024--25.

\textit{Fourth}, we build Layer~0, a real-time pipeline combining GDELT~2.0,
LLM CAMEO event coding, and an AIS tanker equity proxy to deliver sub-monthly
GPR nowcasts between official releases.
Because the pipeline runs on live news data, it can be updated daily---delivering
timely risk signals when geopolitical crises are moving fastest, weeks before
monthly data are available.

\medskip
Section~\ref{sec:lit} reviews related literature.
Section~\ref{sec:data} describes data and GIPI construction.
Section~\ref{sec:decouple} documents the decoupling.
Section~\ref{sec:framework} presents the econometric framework.
Section~\ref{sec:lp} reports local projection results.
Section~\ref{sec:gar} presents GaR results.
Section~\ref{sec:robust} presents three pre-specified robustness checks
confirming the in-sample amplification result.
Section~\ref{sec:layer0} describes Layer~0.
Section~\ref{sec:conclude} concludes.


% ══════════════════════════════════════════════════════════════════════════════
\section{Related Literature}\label{sec:lit}
% ═══════════════════════════════════════════════════════════════════════════════

\noindent\textit{Geopolitical risk and macroeconomic outcomes.}
\citet{caldara2022} construct the GPR index from newspaper text, demonstrating
Granger-causal effects on US investment and industrial production;
critically, their quantile estimates show the 10th-percentile GDP effect
is four times larger than the OLS mean effect, directly motivating a
tail-risk approach.
\citet{caldara2025} document that GPR shocks are robustly inflationary across
44 countries since 1900---a one-standard-deviation shock raises CPI by
0.3--0.5 percentage points over 6--12 months, primarily through energy and
import-price channels.
\citet{pinchetti2024} decompose GPR into an energy component and a broader
macro component, showing that \textit{only} energy-related GPR shocks are
stagflationary; macro-type GPR shocks are deflationary.
This heterogeneity is exactly what GIPI is designed to capture: it measures
whether the energy and pass-through pipes are open when a shock arrives.
\citet{federle2024} use 150 years of data across 60 countries to show that
wars raise CPI by 15 percentage points and cut output by 30\%, with large
geographic spillovers declining with distance---establishing that war-driven
tail risk extends across borders through trade networks \citep{bonadio2021}.
\citet{metiu2025} confirms, in a factor-augmented VAR, that global financial
conditions act as the primary amplification mechanism propagating GPR shocks
internationally---consistent with our inclusion of FCI as a conditioning variable.

\noindent\textit{Policy uncertainty and text-as-data measurement.}
\citet{baker2016}'s EPU index established the newspaper-counting methodology
that Caldara and Iacoviello refined into GPR.
\citet{hassan2019} extend this to firm-level earnings calls, finding that over
90\% of political risk is firm-specific, motivating the use of distributional
methods rather than aggregate averages.
\citet{bloom2009}'s precautionary-investment mechanism explains why output effects
precede any physical disruption, while \citet{kilian2009} provides the foundational
oil-shock decomposition underlying our channel analysis.

\noindent\textit{Growth-at-Risk and quantile methods.}
\citet{adrian2019} establish the foundational GaR result: financial conditions
shift the left tail of the growth distribution procyclically, with lower quantiles
varying up to four times more than the conditional mean.
\citet{adams2021} extend GaR to unemployment and inflation risk distributions,
confirming the framework generalises beyond GDP.
\citet{lopezloria2024} apply GaR to inflation risk directly, finding highly
nonlinear effects of financial conditions on the inflation distribution's tails.
\citet{chavleishvili2024} develop the structural quantile VAR, enabling multivariate
tail forecasts under arbitrary stress scenarios---the natural extension of our approach.
\citet{plagborg2021} prove that local projections and VARs estimate the same
impulse responses under unrestricted lags, providing formal justification for our
quantile LP specification.

\noindent\textit{Supply chains, energy pass-through, and geoeconomic fragmentation.}
\citet{benigno2022} construct the NY Fed GSCPI, showing supply-chain disruptions
transmit inflation globally with a lag of 6--18 months---exactly the horizon at which
our LASSO identifies import prices as the dominant tail predictor.
\citet{baumeister2019} revise upward the share of historical oil price movements
attributable to supply disruptions, reinforcing the energy-channel identification
in our local projections.
\citet{aiyar2023} define geoeconomic fragmentation and estimate global GDP losses
up to 7\% from bloc-based decoupling, with energy trade disruptions causing the
largest price swings---providing the structural context for why the shale supply
revolution dampened the 2023--25 transmission.
\citet{fernandez2024} construct a geopolitical fragmentation index and use local
projections to show fragmentation depresses output and amplifies trade volatility,
complementing GIPI's channel-based measurement.
\citet{caldara2025shortages} construct a shortage index from newspaper text since
1900, showing geopolitical events create shortages that produce output declines
larger relative to their inflation impact---relevant to our finding that IP
responses dominate CPI responses in the LP.
\citet{carriereswallow2023} document state-dependent exchange-rate pass-through,
consistent with our finding that GIPI state (channel openness) determines how
much of a GPR shock reaches consumer prices.

This paper's contribution is to synthesise these literatures for the first time:
we introduce a channel-openness measure (GIPI), embed it in a state-contingent
GaR model via a GPR$\times$GIPI interaction, and provide real-time surveillance
infrastructure to monitor whether the 2025 decoupling persists.
No prior paper has formalised the amplification hypothesis in a quantile tail-risk
framework, identified the specific decoupling mechanism, or built an integrated
Layer~0 nowcasting pipeline.

% ═══════════════════════════════════════════════════════════════════════════════
\section{Data and GIPI Construction}\label{sec:data}
% ═══════════════════════════════════════════════════════════════════════════════

\subsection{Core Monthly Dataset}

The baseline dataset spans January 1985 to December 2025 at monthly frequency.
The GPR index and financial market series are available from 1985;
the regression sample is restricted to January 1997 ($N = 348$ months)
by the availability of the NY~Fed GSCPI, which is required for GIPI
construction.
Industrial production (INDPRO), CPI, core CPI, unemployment, the HY spread
(BAML H0A0), term spread (T10Y2Y), Fed Funds, WTI crude (DCOILWTICO), VIX,
and S\&P~500 total return are from FRED and Yahoo Finance.
The Caldara--Iacoviello GPR index is from \texttt{matteoiacoviello.com}.
Table~\ref{tab:desc} reports key current-period values.

\begin{table}[htbp]\begin{center}
\caption{Key Variables: December 2025 (Most Recent Available)}
\label{tab:desc}
\begin{threeparttable}\small
\begin{tabular}{lccc}
\toprule
Variable & Dec 2025 & Nov 2025 & Sample mean\\
\midrule
GPR Level & ---\tnote{a} & 371.1 & 138\\
GPR $z$-score & --- & 2.65 & 0.0\\
IP YoY growth (\%) & 1.30 & 2.10 & 0.9\\
CPI inflation YoY (\%) & 2.65 & 2.70 & 2.5\\
Core CPI YoY (\%) & 2.65 & 2.60 & 2.4\\
Unemployment (\%) & 4.4 & 4.5 & 5.7\\
Fed Funds rate (\%) & 3.72 & 3.88 & 3.1\\
HY OAS spread (pp) & 2.81 & 2.92 & 5.4\\
VIX & 14.3 & 16.4 & 19.8\\
WTI crude (\$/bbl) & 57.3 & 58.6 & 62.4\\
5Y TIPS breakeven (\%) & 2.26 & 2.29 & 2.1\\
GIPI & $-1.01$ & $-0.83$ & 0.00\\
\bottomrule
\end{tabular}
\begin{tablenotes}[flushleft]\footnotesize
\item[a] GPR is released with a one-month lag.
Sample mean computed over 1985--2025.
\end{tablenotes}
\end{threeparttable}
\end{center}\end{table}

\subsection{Inflation Transmission Channels}

Figure~\ref{fig:channels} displays the six principal inflation transmission
channels over 2023--2025.
The panels show (top row, left to right): CPI YoY (\%; shown for context as the outcome variable), Import Price YoY (\%),
Global Energy YoY (\%); (bottom row): Global Food YoY (\%), 5-Year TIPS Breakeven (\%), and GSCPI (standardised).
The 2022 post-pandemic peaks are visible across all panels: CPI above 8\%,
global energy above $+100\%$ YoY, and food prices above $+40\%$ YoY.
By December 2025, all channels have normalised: energy $-17\%$ YoY,
food $-4.3\%$ YoY, import prices near zero, and GSCPI near its
long-run mean.

\begin{figure}[htbp]\centering
\includegraphics[width=\textwidth]{figures_white/fig01_channels.png}
\caption{Inflation Transmission Channels, 2023--2025.
\textit{Top row, left to right}: CPI YoY (\%; outcome variable, shown for context), Import Price YoY (\%), Global Energy Price YoY (\%).
\textit{Bottom row}: Global Food Price YoY (\%), 5-Year TIPS Breakeven (\%), GSCPI (standardised).
Of the six panels, five enter the GIPI PCA (Import Price, Energy, Food, Breakeven, GSCPI); CPI YoY is excluded as it is the outcome variable.
Annotated end-values as of December 2025.
All GIPI input channels are below their post-pandemic peaks, driving
GIPI to its current reading of $-1.01$ (December 2025).}
\label{fig:channels}
\end{figure}

\subsection{GIPI Construction}

GIPI is the first principal component of five standardised monthly series:
\begin{enumerate}[label=(\arabic*),leftmargin=0.75in,itemsep=2pt,topsep=4pt]
\item NY Fed Global Supply Chain Pressure Index (GSCPI) \citep{benigno2022}
\item Import Price Index YoY (12-month log-difference)
\item IMF Global Energy Price Index YoY (PNRGINDEXM)
\item IMF Global Food Price Index YoY (PFOODINDEXM)
\item $\Delta\text{Breakeven}_{5Y}$: monthly first-difference of T5YIE
\end{enumerate}
Each series is standardised to zero mean and unit variance before PCA.
The sign convention fixes positive loadings on energy, food, and supply-chain
variables so that higher GIPI corresponds to greater inflationary pressure.
GIPI is defined as the PC1 score vector rescaled to unit variance over the
full estimation sample; the resulting standard deviation of 1.524 over
the regression sample reflects sample variation after trimming.\footnote{%
  The raw PC1 score vector has implied SD $=\sqrt{2.715} \approx 1.648$
  from the eigenvalue.
  The reported SD of 1.524 arises from the effective regression sample
  ($N=336$, after trimming for the $h=12$ horizon), where the
  sample moments differ slightly from the full PCA estimation sample
  ($N=348$).}
GIPI spans $N = 348$ months (January 1997 to December 2025);
after trimming for the maximum forecast horizon ($h=12$),
the effective regression sample is $N=336$ observations.

Figure~\ref{fig:gipi} plots the GIPI time series alongside its components.
The 2021--22 episode drives GIPI to its sample maximum of approximately $+5.6$,
reflecting near-simultaneous surges in all five channels.
The COVID dip in mid-2020 (GIPI below $-2$) is clearly visible.
The current reading of $-1.01$ reflects simultaneous
below-average readings on all five channels.

\begin{figure}[htbp]\centering
\includegraphics[width=\textwidth]{figures_white/fig02_gipi.png}
\caption{GIPI: Geopolitical Inflation Pressure Index (Detail: 2023--2025; full sample 1997--2025).
\textit{Gold solid}: GIPI (PC1 of five components).
\textit{Blue fill}: negative GIPI (channels subdued).
\textit{Gold fill}: positive GIPI (channels elevated).
\textit{Orange dashed}: elevated threshold ($+3.1$).
\textit{Red dashed}: crisis threshold ($+5.2$).
Note: thresholds are set on the raw PC1 score scale, calibrated to the
approximate GIPI values observed at the 85th ($+3.1$) and 97th ($+5.2$)
sample percentiles of the full 1997--2025 distribution; the GPR regime
thresholds of $z=1.5$ and $z=2.5$ are applied to the GPR $z$-score,
not directly to GIPI.
Overlaid dashed lines: Import Price YoY (teal, $\div$1), Energy YoY (orange, $\div$20),
Food YoY (green, $\div$10); divisors chosen for visual scale only,
to allow directional comparison on a common axis.
Current reading (December 2025): $\mathbf{-1.01}$, subdued regime.
Post-pandemic peak (June 2022): $\approx+5.6$, crisis regime.}
\label{fig:gipi}
\end{figure}

\subsection{GIPI Validation}

Table~\ref{tab:pca} reports the PCA output for the full 1997--2025 sample.
The first principal component explains 54.3\% of common variance in the
five transmission-channel series, substantially more than the second
component (20.1\%), supporting its interpretation as a single
``inflationary pressure'' factor.
All five loadings carry the expected positive sign after the sign
normalisation: energy prices load most heavily ($0.531$), followed by
food prices ($0.494$) and GSCPI ($0.461$), with import prices ($0.380$)
and breakeven changes ($0.337$) carrying somewhat lower weight.
The dominance of energy and food price channels is consistent with
the commodity-transmission literature \citep{conflitti2019}.

\begin{table}[htbp]\begin{center}
\caption{GIPI Principal Component Analysis: Loadings and Variance Explained}
\label{tab:pca}
\begin{threeparttable}\small\setlength{\tabcolsep}{8pt}
\begin{tabular}{@{}l r r r@{}}
\toprule
Series & PC1 Loading & PC2 Loading & Communality\\
\midrule
Global Energy YoY       & $0.531$ & $-0.182$ & $0.315$\\
Global Food YoY         & $0.494$ & $-0.238$ & $0.301$\\
GSCPI                   & $0.461$ & $+0.421$ & $0.390$\\
Import Price YoY        & $0.380$ & $+0.509$ & $0.403$\\
$\Delta$Breakeven$_{5Y}$& $0.337$ & $+0.681$ & $0.578$\\
\midrule
Variance explained      & 54.3\%  & 20.1\%   & ---\\
Cumulative              & 54.3\%  & 74.4\%   & ---\\
\bottomrule
\end{tabular}
\begin{tablenotes}[flushleft]\footnotesize
\item \textit{Notes:} PCA estimated on standardised series, January 1997--December 2025
($N=348$). Sign convention: positive loadings correspond to inflationary pressure.
PC1 eigenvalue $= 2.715$; PC2 eigenvalue $= 1.005$.
Communality $=$ share of each variable's variance captured by PC1 and PC2 jointly.
The 54.3\% variance share for PC1 substantially exceeds the $1/K = 20\%$ threshold,
supporting retention as a single dominant factor.
The breakeven change ($\Delta$Breakeven$_{5Y}$) loads primarily on PC2 (a
``financial conditions'' factor), not PC1, reflecting its partially independent
information content from the real commodity channels.
This implies that the Fed credibility mechanism (Section~\ref{sec:decouple},
mechanism iii) operates partly through a second factor that GIPI's PC1
construction underweights; a PC1$+$PC2 specification (discussed as an
extension in Section~\ref{sec:conclude}) would better capture this channel.
\end{tablenotes}
\end{threeparttable}
\end{center}\end{table}

The loading structure is stable across sub-samples.
Estimated on 1997--2014 (pre-shale, pre-pandemic), PC1 still explains 51.8\%
of variance and the energy loading remains the largest ($0.523$), confirming
that the 2021--22 pandemic surge does not mechanically dominate the factor.
The 2023--25 CALM readings reflect genuine below-average realisations on
all five inputs, not a statistical artefact of the full-sample PCA.

% ═══════════════════════════════════════════════════════════════════════════════
\section{The GPR--GIPI Decoupling}\label{sec:decouple}
% ═══════════════════════════════════════════════════════════════════════════════

Figure~\ref{fig:gpr} plots the GPR level and $z$-score from 2023--2025 with
regime shading.
The crisis months of 2025 are clearly visible: April 2025 reached GPR $= 623.4$
($z = 5.63$), the highest sustained reading since the post-9/11
period.\footnote{%
  The $z$-score is standardised against the 1997--2025 regression sample
  (mean $= 138$, SD $= 86$).
  All GPR data in this paper are sourced from the FRED series
  GEPUCURRENT (downloaded March~2026).
  The Caldara--Iacoviello index undergoes periodic revisions;
  later vintages of the GPR index available on the authors' website
  may differ in both level and ranking of historical episodes.
  Comparisons of raw GPR levels across distant episodes should be
  interpreted with this caveat in mind.}

\begin{figure}[htbp]\centering
\includegraphics[width=\textwidth]{figures_white/fig03_gpr.png}
\caption{Geopolitical Risk Index Monitor, 2023--2025.
\textit{Orange solid}: GPR level (left axis, scale 0--700).
\textit{Blue dashed}: GPR $z$-score (right axis).
\textit{Background shading}: green = calm regime, amber = elevated, red = crisis.
Dotted horizontal lines mark the elevated ($z=1.5$) and crisis ($z=2.5$) thresholds.
November 2025 reading annotated: GPR $= 371.1$, $z = 2.65$ (crisis band).
Note: GPR released with one-month lag; December 2025 not yet available.}
\label{fig:gpr}
\end{figure}

Figure~\ref{fig:decoupling} documents the decoupling directly.
The left panel scatters GPR $z$-score against GIPI, coloured by regime:
all crisis and elevated observations (orange and red) cluster in the
lower half of the chart (GIPI $< 0$), far from the upper-right quadrant
that characterised the 2022 episode.
The in-sample correlation is $\rho = 0.131$ over the 2023--25 window.
The right panel displays the two series as a dual-axis time series,
confirming the sustained divergence.

\begin{figure}[htbp]\centering
\includegraphics[width=\textwidth]{figures_white/fig04_decoupling.png}
\caption{The GPR--GIPI Decoupling, 2023--2025.
\textit{Left panel}: Scatter of GPR $z$-score vs GIPI coloured by regime
(green = calm, orange = elevated, red = crisis).
Vertical dotted lines at $z=1.5$ and $z=2.5$ mark regime boundaries.
Horizontal dashed line at GIPI $= 0$.
All crisis and elevated observations show negative GIPI; $\rho = 0.131$.
\textit{Right panel}: Dual-axis time series of GPR $z$-score (orange, left axis)
and GIPI (blue dashed, right axis) showing the sustained divergence.
Orange dotted: crisis threshold ($z=2.5$). Blue dashed at 0: GIPI neutral level.}
\label{fig:decoupling}
\end{figure}

Table~\ref{tab:decouple} lists every crisis and elevated episode with
contemporaneous GIPI values.
In all seven crisis months, GIPI ranged from $-1.04$ to $-1.59$.
The April 2025 observation---GPR $z = 5.63$ paired with GIPI $= -1.54$,
energy $-15.4\%$ YoY---is the paper's sharpest identification point.

\begin{table}[htbp]\begin{center}
\caption{Crisis and Elevated Episodes: GPR vs GIPI, 2023--2025}
\label{tab:decouple}
\begin{threeparttable}\small
\begin{tabular}{lcccccl}
\toprule
Date & GPR & $z$ & Regime & GIPI & Energy YoY & Imp.\ $P$ YoY\\
\midrule
2023-03 & 313.7 & 1.97 & Elevated & $-2.45$ & $-43.9\%$ & $-4.7\%$\\
2025-01 & 309.1 & 1.91 & Elevated & $-0.50$ & $+8.6\%$ & $+1.7\%$\\
2025-02 & 327.7 & 2.13 & Elevated & $-0.75$ & $+5.7\%$ & $+1.7\%$\\
\rowcolor{yellow!15}
2025-03 & 468.8 & 3.80 & \textbf{Crisis} & $-1.17$ & $-3.9\%$ & $+0.8\%$\\
\rowcolor{yellow!15}
2025-04 & \textbf{623.4} & \textbf{5.63} & \textbf{Crisis} & $-1.54$ & $-15.4\%$ & $+0.0\%$\\
\rowcolor{yellow!15}
2025-05 & 511.7 & 4.31 & \textbf{Crisis} & $-1.59$ & $-15.7\%$ & $-0.4\%$\\
\rowcolor{yellow!15}
2025-06 & 368.6 & 2.62 & \textbf{Crisis} & $-1.26$ & $-10.2\%$ & $-0.6\%$\\
\rowcolor{yellow!15}
2025-07 & 406.9 & 3.07 & \textbf{Crisis} & $-1.04$ & $-11.2\%$ & $-0.4\%$\\
2025-08 & 343.3 & 2.32 & Elevated & $-1.21$ & $-12.9\%$ & $-0.3\%$\\
2025-09 & 308.5 & 1.91 & Elevated & $-0.94$ & $-7.4\%$ & $-0.1\%$\\
\rowcolor{yellow!15}
2025-10 & 404.9 & 3.05 & \textbf{Crisis} & $-1.32$ & $-14.6\%$ & $-0.2\%$\\
\rowcolor{yellow!15}
2025-11 & 371.1 & 2.65 & \textbf{Crisis} & $-1.18$ & $-13.9\%$ & $-0.1\%$\\
\bottomrule
\end{tabular}
\begin{tablenotes}[flushleft]\footnotesize
\item \textit{Notes:} Highlighted rows = crisis regime.
Every crisis month has negative GIPI.
The April 2025 observation (GPR $z = 5.63$, the 2025 peak) coincides with
GIPI $= -1.54$: energy prices $-15.4\%$ YoY, driven by record US production
and IEA reserve releases, preventing a global supply shortfall despite
intense Middle East conflict.
\end{tablenotes}
\end{threeparttable}
\end{center}\end{table}

Three structural mechanisms explain the decoupling.
\textit{(i) US shale supply response:} US crude production reached
record highs of 13.5\,mb/d in late 2025, eliminating the global
supply shortfall that characterised pre-shale Middle East crises.
\textit{(ii) IEA reserve coordination:} Coordinated strategic
reserve releases in 2023--24 capped Brent spikes within days of
escalation announcements.
\textit{(iii) Fed credibility:} With the Fed Funds rate at 3.72\%
and expectations well-anchored, TIPS breakevens did not reprice
the de-anchoring scenario, keeping the $\Delta\text{Breakeven}_{5Y}$
component of GIPI near zero.

% ═══════════════════════════════════════════════════════════════════════════════
\section{Econometric Framework}\label{sec:framework}
% ═══════════════════════════════════════════════════════════════════════════════

\noindent\textbf{Notation.}
Two GPR constructs appear in the specifications below.
\textit{GPR\_shock}$_t$ denotes the standardised monthly change
$\Delta\text{GPR}_t / \sigma_{\Delta\text{GPR}}$ and is used in the
local projections (Equation~\ref{eq:lp}), where the impulse is a
one-time monthly innovation.
\textit{GPR}$_t$ denotes the standardised level
$(\text{GPR}_t - \bar{\text{GPR}}) / \sigma_{\text{GPR}}$ (i.e.\ the
$z$-score) and is used in the GaR specifications
(Equations~\ref{eq:base}--\ref{eq:enhanced}), where the conditioning
variable is the \textit{state} of geopolitical risk.

\subsection{Financial Conditions Index}

The FCI is the first principal component of five standardised monthly series:
log VIX, the HY OAS spread, the 10Y--2Y term spread, $\Delta$Fed~Funds,
and the S\&P~500 monthly return.
Higher FCI values indicate tighter financial conditions.

\subsection{Local Projections}

For outcomes $y \in \{\text{IP YoY},\ \text{CPI},\ \text{unemployment}\}$
and horizons $h = 1,\ldots,24$ \citep{jorda2005}:
\begin{equation}
y_{t+h} - y_{t-1} = \alpha_h + \beta_h\,\text{GPR\_shock}_t
+ \boldsymbol{\Gamma}_h\,\mathbf{X}_t
+ \sum_{k=1}^{4}\phi_{h,k}\,y_{t-k} + \varepsilon_{t+h}.
\label{eq:lp}
\end{equation}
GPR\_shock$_t$ is the standardised monthly GPR change.
Standard errors follow \citet{newey1987}, bandwidth $h+1$.

\subsection{Growth-at-Risk: Baseline}

\begin{equation}
Q_\tau(y_{t+h} \mid \mathbf{x}_t)
= \alpha_\tau + \beta_\tau\,\text{GPR}_t
+ \gamma_\tau\,\text{FCI}_t + \eta_\tau\,y_t,
\label{eq:base}
\end{equation}
for $\tau \in \{0.05, 0.10, 0.25, 0.50, 0.75, 0.90, 0.95\}$,
$h \in \{3, 6, 12\}$.

\subsection{Enhanced Specification}

\begin{equation}
Q_\tau(y_{t+h} \mid \mathbf{x}_t)
= \alpha_\tau + \beta_\tau\,\text{GPR}_t + \gamma_\tau\,\text{FCI}_t
+ \delta_\tau\,\text{GIPI}_t
+ \zeta_\tau\,(\text{GPR}_t \times \text{GIPI}_t)
+ \eta_\tau\,y_t.
\label{eq:enhanced}
\end{equation}
Both GPR and GIPI are standardised before the interaction.
The prediction is $\hat{\zeta}_\tau < 0$ at $\tau \leq 0.10$:
when GIPI is high, a GPR shock inflicts greater left-tail damage.
With GIPI $= -1.01$, the interaction term instead attenuates tail risk.
Forecast accuracy: mean pinball loss $\mathcal{L}_\tau$ and
Koenker--Machado pseudo-$R^2_\tau$ \citep{koenker1978,koenker1999}.

\subsection{Three Pre-Specified Robustness Checks}

\textit{Check 1 (Orthogonalisation).}---Replace GIPI with its
OLS residual from a regression on GPR\_shock (GIPI$^\perp$).
Significance of $\hat{\delta}^\perp_\tau < 0$ at $\tau \leq 0.10$
confirms GIPI's independent left-tail content.

\textit{Check 2 (LASSO-QR).}---Expand to all five individual
GIPI components; apply $L_1$-penalised quantile regression
(5-fold CV at $\tau = 0.10$).
Surviving variables identify the robust component-level predictors.

\textit{Check 3 (Scoring Rules).}---Compare $\mathcal{L}_\tau$ and
$R^2_\tau$ between equations~\eqref{eq:base} and~\eqref{eq:enhanced}
at all 21 ($\tau$, $h$) cells.

% ═══════════════════════════════════════════════════════════════════════════════
\section{Local Projection Results}\label{sec:lp}
% ═══════════════════════════════════════════════════════════════════════════════

Figure~\ref{fig:lp_irf} displays the LP impulse response functions
for three outcome variables at horizons 0--24 months.

\begin{figure}[htbp]\centering
\includegraphics[width=\textwidth]{figures_white/fig05_lp_irf.png}
\caption{Local Projection Impulse Responses: One Standard-Deviation GPR Shock.
\textit{Left panel}: Response of IP YoY growth (percentage points).
The response is monotonically negative, becoming statistically significant
(68\% CI excludes zero) from approximately $h=8$ months, reaching
$-0.9$\,pp at $h=18$.
\textit{Centre panel}: Response of CPI inflation YoY.
The CPI response is muted and drifts slightly negative over
the forecast horizon---near zero for $h \leq 6$ and drifting
toward $-0.18$\,pp at $h \approx 15$---reflecting the
paper's central decoupling thesis: under current conditions,
GPR shocks do not transmit to consumer prices through the energy
and import-price channels when GIPI is subdued.
\textit{Right panel}: Response of unemployment rate.
Small positive effects, barely significant at conventional levels,
consistent with the precautionary investment delay channel \citep{bloom2009}
rather than immediate layoffs.
All panels: dark shaded band = 68\% bootstrap interval; light band = 90\%.
Horizontal dashed line at zero.}
\label{fig:lp_irf}
\end{figure}

The IP response is the paper's primary LP finding.
The monotonic deepening through $h=18$---with the 68\% confidence band
excluding zero from $h=8$ onward---confirms that GPR shocks create
persistent rather than transitory output effects.
The CPI response drifts slightly negative over the horizon,
consistent with the LASSO-QR finding in Section~\ref{sec:robust}:
the current low-GIPI environment suppresses the import-price
pass-through that would otherwise transmit GPR shocks to headline inflation.
The unemployment response's marginal significance at short horizons is
consistent with \citet{bloom2009}: firms reduce investment and hiring
plans before any physical disruption affects revenues, creating a
precautionary channel distinct from the direct output contraction.

% ═══════════════════════════════════════════════════════════════════════════════
\section{Growth-at-Risk Results}\label{sec:gar}
% ═══════════════════════════════════════════════════════════════════════════════

GIPI contains energy prices, food prices, supply-chain pressure,
and breakevens --- variables that respond both to GPR shocks (supply
channel) and to US IP growth itself (demand channel).
A contemporaneous specification $\text{GPR}_t \times \text{GIPI}_t$ therefore
risks simultaneity bias: a positive IP shock raises energy demand and thus
GIPI, conflating the demand-driven channel with the geopolitical-amplification
channel.

We address this in three ways.
\textit{First}, our baseline uses \textbf{GIPI$_{t-1}$} in the interaction:
pre-shock channel openness is predetermined with respect to $y_{t+h}$
and cannot be driven by future output realisations.
\textit{Second}, the local projection design (Equation~1) uses current GPR
to predict \textit{future} IP growth at $h = 3, 6, 12$ months forward,
substantially reducing contemporaneous simultaneity.
\textit{Third}, Check~1 (Table~\ref{tab:orth}) orthogonalises GIPI on the
GPR shock, removing any linear dependence; the interaction coefficient
remains significant, confirming the result is not an artefact of
GPR$\to$GIPI pass-through.

The direction of residual bias is conservative.
Demand-driven GIPI increases (boom $\to$ high oil demand $\to$ high GIPI)
occur precisely when IP growth is high --- the \textit{right} tail, not the
left tail we care about.
Any remaining endogeneity therefore attenuates the amplification coefficient
$\hat\zeta_\tau$ at $\tau \leq 0.10$, making our estimates
conservative lower bounds on the true left-tail amplification effect
(Section~\ref{sec:iv}).

\subsection{Identification Scope}
\label{sec:iv}

\noindent\textbf{What the paper identifies.}
All results should be read as \textit{conditional in-sample associations}:
given that inflation transmission channels are in state GIPI$_{t-1}$,
the historical co-movement between GPR$_z$ and the left tail of IP growth
is characterised by the quantile regression coefficients reported above.
We do not claim to recover structural causal effects of geopolitical risk
on macroeconomic outcomes.

Supply shocks can simultaneously raise both GPR and GIPI --- as in
the Russia--Ukraine episode of early 2022, when the initial invasion
lifted the GPR index and simultaneously triggered food and energy price
shocks that raised GIPI --- creating a positive co-movement between the
regressor and the moderator that GIPI$_{t-1}$ predetermination does not
fully neutralise.
A formal instrument for GIPI, such as the exogenous oil supply news
series of \citet{baumeister2019} or OPEC spare-capacity shocks
\citep{caldara2022}, would strengthen causal interpretation and is a
natural extension of this work.
The three strategies we employ (lagged GIPI, forward projection design,
orthogonalisation in Check~1) reduce but do not eliminate the
identification challenge; the conservative-bias argument means our
left-tail amplification estimates are lower bounds under reasonable
assumptions about the sign of any residual endogeneity.

\subsection{Baseline vs.\ Enhanced Comparison}

Table~\ref{tab:garfull} compares baseline and enhanced specifications
across all six estimation rows.

\begin{table}[htbp]\begin{center}
\caption{Growth-at-Risk: Baseline vs.\ Enhanced Model (December 2025 Nowcast)}
\label{tab:garfull}
\begin{threeparttable}\small
\begin{tabular}{llccccccc}
\toprule
$h$ & Spec & Median & GaR$_5$ & GaR$_{25}$ & $\Pr(y{<}0)$ & $\Pr(\text{rec.})$ & Skew\\
&& (\%) & (\%) & (\%) & & &\\
\midrule
\multirow{3}{*}{3}
& Baseline & $+1.12$ & $-1.87$ & $+0.21$ & 22.6\% & 1.8\% & $+1.02$\\
& Enhanced & $+1.30$ & $-1.25$ & $+0.48$ & 16.7\% & 0.7\% & $+0.53$\\
& \textit{$\Delta$} & \textit{$+0.18$} & \textit{$+0.62$} & \textit{$+0.27$} &
  \textit{$-5.9$pp} & \textit{$-1.1$pp} & \textit{$-0.49$}\\
\midrule
\multirow{3}{*}{6}
& Baseline & $+0.93$ & $-3.20$ & $-0.22$ & 31.4\% & 6.4\% & $+0.68$\\
& \textbf{Enhanced} & $\mathbf{+1.30}$ & $-3.16$ & $-0.42$ & \textbf{22.4\%} &
  \textbf{2.7\%} & $+0.27$\\
& \textit{$\Delta$} & \textit{$+0.37$} & \textit{$+0.04$} & \textit{$-0.20$} &
  \textit{$-9.0$pp} & \textit{$-3.7$pp} & \textit{$-0.41$}\\
\midrule
\multirow{3}{*}{12}
& Baseline & $-0.27$ & $-6.19$ & $-1.47$ & 54.4\% & 24.4\% & $+1.39$\\
& \textbf{Enhanced} & $\mathbf{+0.28}$ & $-7.53$ & $-1.98$ & \textbf{46.1\%} &
  \textbf{21.3\%} & $+0.66$\\
& \textit{$\Delta$} & \textit{$+0.55$} & \textit{$-1.34$} & \textit{$-0.51$} &
  \textit{$-8.3$pp} & \textit{$-3.1$pp} & \textit{$-0.73$}\\
\bottomrule
\end{tabular}
\begin{tablenotes}[flushleft]\footnotesize
\item \textit{Notes:}
Baseline: GPR$+$FCI$+y_t$.
Enhanced: Baseline$+$GIPI$_{t-1}+$(GPR$\times$GIPI$_{t-1}$).
$\Pr(\text{rec.}) = \Pr(y_{t+h}<-2\%)$, where $-2\%$ YoY IP growth
serves as a proxy for severe contraction broadly consistent with
NBER recession episodes in our sample; we use this label for
brevity, acknowledging that IP growth alone does not fully determine
NBER dating.
Skew $=(Q_{75}-Q_{50})/(Q_{50}-Q_{25})$.
Current conditions: GIPI~$=-1.01$ (December 2025), GPR $z=2.65$ (November 2025, 1-month lag).
$N=336$, 1997M01--2025M12.
\textbf{Bold}: enhanced specification rows at $h=6$ and $h=12$.
\end{tablenotes}
\end{threeparttable}
\end{center}\end{table}

\subsection{Fan Charts}

Figures~\ref{fig:fan3}--\ref{fig:fan12} display the conditional predictive
densities from the enhanced specification at $h \in \{3, 6, 12\}$ months.
For each horizon the density is shaded by outcome zone (recession,
contraction, expansion), with vertical lines marking the conditional
median and GaR$_5$ (5th quantile).

\begin{figure}[p]\centering
\includegraphics[width=\textwidth]{figures_white/fig06_fan_h3.png}
\caption{Conditional Predictive Density, $h=3$ months, Enhanced Model.
\textit{Shaded regions}: recession (IP $< -2\%$, red), contraction ($-2\%$ to $0\%$, orange),
expansion ($> 0\%$, green).
\textit{Vertical lines}: conditional median (blue dashed) and GaR$_5$ (red dashed).
December 2025 nowcast: GaR$_5 = -1.2\%$, Median $= +1.3\%$,
$\Pr(\text{IP}<0) = 16.7\%$, $\Pr(\text{rec.}) = 0.7\%$,
conditional skewness $= +0.53$.
At 3 months, tail risk is tightly bounded under subdued GIPI.}
\label{fig:fan3}
\end{figure}

\begin{figure}[p]\centering
\includegraphics[width=\textwidth]{figures_white/fig06_fan_h6.png}
\caption{Conditional Predictive Density, $h=6$ months, Enhanced Model.
Same layout as Figure~\ref{fig:fan3}.
December 2025 nowcast: GaR$_5 = -3.2\%$, Median $= +1.3\%$,
$\Pr(\text{IP}<0) = 22.4\%$, $\Pr(\text{rec.}) = 2.7\%$,
conditional skewness $= +0.27$.
This is the paper's headline six-month risk assessment:
the central case is benign (median $+1.3\%$); GaR$_5$ captures
the scenario in which inflation channels re-engage with current
geopolitical risk despite currently subdued GIPI conditions.
Note: the GIPI enhancement at $h=6$ is concentrated in the median
and recession probability rather than GaR$_5$ itself ($-3.2\%$
vs.\ $-3.2\%$ baseline), reflecting the intermediate horizon at which
channel pass-through has not yet fully accumulated.}
\label{fig:fan6}
\end{figure}

\begin{figure}[p]\centering
\includegraphics[width=\textwidth]{figures_white/fig06_fan_h12.png}
\caption{Conditional Predictive Density, $h=12$ months, Enhanced Model.
December 2025 nowcast: GaR$_5 = -7.5\%$, Median $= +0.3\%$,
$\Pr(\text{IP}<0) = 46.1\%$, $\Pr(\text{rec.}) = 21.3\%$,
conditional skewness $= +0.66$.
At the 12-month horizon, the fan widens substantially,
reflecting accumulated path uncertainty in both GPR and GIPI.
The enhanced model lifts the median ($+0.3\%$ vs $-0.3\%$ baseline)
and reduces the recession probability (21.3\% vs 24.4\%).}
\label{fig:fan12}
\end{figure}

\FloatBarrier

\subsection{Conditional Predictive Densities}

Figure~\ref{fig:densities} presents the three conditional predictive
densities side-by-side for $h \in \{3, 6, 12\}$.
Each panel shows the asymmetric skew-normal approximation to the enhanced
model's conditional distribution, with labelled quantile lines at the
5th, 25th, 50th, 75th, and 95th percentiles.

\begin{figure}[htbp]\centering
\includegraphics[width=\textwidth]{figures_white/fig07_densities.png}
\caption{Conditional Predictive Densities, Enhanced Model, December 2025 Nowcast.
\textit{Three panels}: $h=3$ months (left), $h=6$ months (centre), $h=12$ months (right).
\textit{Blue fill}: full density; darker blue fill: interquartile range (25th--75th).
\textit{Vertical lines}: 5th (red dashed), 25th (blue dotted), 50th/median (blue solid),
75th (blue dotted), 95th (red dashed) quantiles, annotated with their percentage-point values.
\textit{Grey dashed}: $-2\%$ recession threshold.
Titles display the horizon-specific recession probability and conditional skewness.
Distributions widen with horizon as captured by GaR$_5$
($-1.3\% \to -3.2\% \to -7.5\%$), while positive Bowley skewness
($+0.53 \to +0.27 \to +0.66$) reflects upside potential under
currently subdued GIPI conditions ($-1.01$).}
\label{fig:densities}
\end{figure}

\FloatBarrier

% ═══════════════════════════════════════════════════════════════════════════════
\section{Robustness Checks}\label{sec:robust}
% ═══════════════════════════════════════════════════════════════════════════════

Figure~\ref{fig:robust} provides a visual summary of all three checks.

\begin{figure}[htbp]\centering
\includegraphics[width=\textwidth]{figures_white/fig08_robustness.png}
\caption{Robustness Summary: Three Pre-Specified Checks.
\textit{Left panel (Check 1)}: Orthogonalised GIPI$^\perp$ coefficient
$\hat{\delta}^\perp_\tau$ by quantile $\tau$ and horizon (blue = $h=3$m,
teal = $h=6$m, red = $h=12$m). Bars below zero at $\tau \leq 0.25$ confirm
GIPI carries independent left-tail information at all horizons.
Positive bars at $\tau \geq 0.75$ reflect the energy-sector output windfall
at high GIPI.
\textit{Centre panel (Check 3)}: Bubble chart of pinball improvement (\%)
for all 21 quantile-horizon cells. Bubble area proportional to improvement;
colour scale from yellow (small) to red (large). Headline result:
$+14.7\%$ at ($h=12$, $\tau=5\%$).
\textit{Right panel (Check 3)}: Pseudo-$R^2$ by quantile, baseline (dashed)
vs enhanced (solid) for all three horizons. The enhanced model weakly
dominates at every cell.}
\label{fig:robust}
\end{figure}

Tables~\ref{tab:orth}--\ref{tab:lasso} and \ref{tab:pinball} provide the full numerical results.

\begin{table}[htbp]\begin{center}
\caption{Check 1: Orthogonalisation --- GIPI$^\perp$ Coefficient $\hat{\delta}^\perp_\tau$}
\label{tab:orth}
\begin{threeparttable}\small\setlength{\tabcolsep}{8pt}
\begin{tabular}{@{}l *{6}{c}@{}}
\toprule
 & \multicolumn{6}{c}{Quantile $\tau$}\\
\cmidrule(l){2-7}
Horizon & 0.05 & 0.10 & 0.25 & 0.50 & 0.75 & 0.95\\
\midrule
$h=3$  & $-0.618^{**}$ & $-0.107^{*}$   & $-0.118^{*}$   & $+0.209^{**}$ & $+0.377^{***}$ & $+1.722^{***}$\\[2pt]
$h=6$  & $-1.106^{***}$ & $-0.884^{***}$ & $-0.071^{**}$ & $+0.120^{*}$  & $+0.030$       & $+0.783^{**}$\\[2pt]
$h=12$ & $-2.104^{***}$ & $-2.170^{***}$ & $+0.442^{*}$   & $-0.054^{*}$ & $+0.188$       & $-0.522^{*}$\\
\bottomrule
\end{tabular}
\begin{tablenotes}[flushleft]\footnotesize
\item \textit{Notes:} GIPI$^\perp$ = residual from OLS projection of GIPI on GPR\_shock.
16 of 18 cells are significant at the 10\% level or better;
the two exceptions ($h{=}6$, $\tau{=}0.75$ and $h{=}12$, $\tau{=}0.75$)
are economically small and in the right tail where the amplification
prediction does not apply.
Newey--West SE with bandwidth $h{+}1$.
$^{*}p<0.10$;\; $^{**}p<0.05$;\; $^{***}p<0.01$.
The $h{=}12$ lower-tail coefficients ($-2.10$, $-2.17$) represent the
strongest evidence for independent GIPI information content.
$\tau = 0.90$ is omitted from this table for space; the corresponding
enhanced-model scoring-rule improvements are reported in
Table~\ref{tab:pinball}.
\end{tablenotes}
\end{threeparttable}
\end{center}\end{table}

\begin{table}[htbp]\begin{center}
\caption{Check 2: LASSO-QR Surviving Predictors and Pinball Improvement}
\label{tab:lasso}
\begin{threeparttable}\small\setlength{\tabcolsep}{7pt}
\begin{tabular}{@{}cc>{\raggedright\arraybackslash}p{6.2cm}r@{}}
\toprule
$h$ & $\tau$ & Surviving predictors & Impr.(\%)\\
\midrule
\multirow{7}{*}{3}
 & 0.05 & FCI, ip\_yoy                                          & $+2.8$\\
 & 0.10 & FCI, ip\_yoy                                          & $+0.3$\\
 & 0.25 & FCI, ip\_yoy, $\Delta$bkeven$_{5Y}$                   & $+0.3$\\
 & 0.50 & GPR$_z$, ip\_yoy, import\_price\_yoy                  & $+0.7$\\
 & 0.75 & GPR$_z$, ip\_yoy, import\_price\_yoy                  & $+1.6$\\
 & 0.90 & ip\_yoy, import\_price\_yoy                           & $+2.8$\\
 & 0.95 & global\_energy\_yoy                                   & $+8.1$\\
\midrule
\multirow{7}{*}{6}
 & 0.05 & FCI, ip\_yoy                                          & $+7.3$\\
 & 0.10 & FCI, ip\_yoy                                          & $+2.0$\\
 & 0.25 & FCI, ip\_yoy                                          & $+0.2$\\
 & 0.50 & GPR$_z$, ip\_yoy, import\_price\_yoy, food\_yoy       & $+0.2$\\
 & 0.75 & GPR$_z$, ip\_yoy, import\_price\_yoy                  & $+0.1$\\
 & 0.90 & (none survive shrinkage)                              & $+0.3$\\
 & 0.95 & (none survive shrinkage)                              & $+2.2$\\
\midrule
\multirow{7}{*}{12}
 & \cellcolor{yellow!25}0.05 & \cellcolor{yellow!25}\textbf{FCI, import\_price\_yoy}  & \cellcolor{yellow!25}$\mathbf{+14.7}$\\
 & \cellcolor{yellow!25}0.10 & \cellcolor{yellow!25}\textbf{FCI, import\_price\_yoy}  & \cellcolor{yellow!25}$\mathbf{+11.4}$\\
 & 0.25 & FCI, global\_energy\_yoy, $\Delta$bkeven$_{5Y}$       & $+0.9$\\
 & 0.50 & GPR$_z$                                               & $+0.1$\\
 & 0.75 & FCI, ip\_yoy, food\_yoy, $\Delta$bkeven$_{5Y}$        & $+0.2$\\
 & 0.90 & ip\_yoy                                               & $+1.7$\\
 & 0.95 & ip\_yoy                                               & $+5.9$\\
\bottomrule
\end{tabular}
\begin{tablenotes}[flushleft]\footnotesize
\item \textit{Notes:}
$L_1$-penalised quantile regression; 5-fold CV tuning at $\tau=0.10$ per horizon.
Predictor pool: GPR$_z$, FCI, five GIPI components, lagged ip\_yoy.
\textbf{Highlighted}: headline $h=12$ lower-tail results where import prices displace GPR,
consistent with a 9--18-month pass-through lag \citep{conflitti2019}.
\end{tablenotes}
\end{threeparttable}
\end{center}\end{table}

\begin{table}[htbp]\begin{center}
\caption{Check 3: Complete Scoring-Rule Results, Baseline vs.\ Enhanced}
\label{tab:pinball}
\begin{threeparttable}\small\setlength{\tabcolsep}{5.5pt}
\begin{tabular}{@{}cc *{3}{r} @{\hspace{8pt}} *{3}{r}@{}}
\toprule
 &  & \multicolumn{3}{c}{Pinball Loss $\mathcal{L}_\tau$}
    & \multicolumn{3}{c}{Pseudo-$R^2_\tau$}\\
\cmidrule(lr){3-5}\cmidrule(l){6-8}
$h$ & $\tau$ & Base & Enh. & $\Delta$(\%) & Base & Enh. & $\Delta R^2$\\
\midrule
\multirow{7}{*}{3}
 & 0.05 & 0.279 & 0.271 & $+2.8$ & 0.573 & 0.585 & $+0.012$\\
 & 0.10 & 0.398 & 0.397 & $+0.3$ & 0.579 & 0.581 & $+0.002$\\
 & 0.25 & 0.623 & 0.621 & $+0.3$ & 0.545 & 0.546 & $+0.001$\\
 & 0.50 & 0.766 & 0.760 & $+0.7$ & 0.454 & 0.458 & $+0.004$\\
 & 0.75 & 0.622 & 0.613 & $+1.6$ & 0.370 & 0.379 & $+0.010$\\
 & 0.90 & 0.400 & 0.389 & $+2.8$ & 0.272 & 0.293 & $+0.020$\\
 & 0.95 & 0.288 & 0.265 & $+8.1$ & 0.182 & 0.248 & $+0.067$\\
\midrule
\multirow{7}{*}{6}
 & 0.05 & 0.402 & 0.373 & $+7.3$ & 0.384 & 0.429 & $+0.045$\\
 & 0.10 & 0.593 & 0.582 & $+2.0$ & 0.372 & 0.385 & $+0.013$\\
 & 0.25 & 0.954 & 0.952 & $+0.2$ & 0.300 & 0.302 & $+0.002$\\
 & 0.50 & 1.090 & 1.088 & $+0.2$ & 0.217 & 0.218 & $+0.001$\\
 & 0.75 & 0.823 & 0.822 & $+0.1$ & 0.153 & 0.154 & $+0.001$\\
 & 0.90 & 0.503 & 0.502 & $+0.3$ & 0.059 & 0.062 & $+0.002$\\
 & 0.95 & 0.334 & 0.326 & $+2.2$ & 0.029 & 0.050 & $+0.022$\\
\midrule
\multirow{7}{*}{12}
 & \cellcolor{yellow!25}0.05 & \cellcolor{yellow!25}0.559 & \cellcolor{yellow!25}0.477 & \cellcolor{yellow!25}$\mathbf{+14.7}$ & \cellcolor{yellow!25}0.146 & \cellcolor{yellow!25}0.271 & \cellcolor{yellow!25}$\mathbf{+0.125}$\\
 & \cellcolor{yellow!25}0.10 & \cellcolor{yellow!25}0.886 & \cellcolor{yellow!25}0.785 & \cellcolor{yellow!25}$\mathbf{+11.4}$ & \cellcolor{yellow!25}0.065 & \cellcolor{yellow!25}0.172 & \cellcolor{yellow!25}$\mathbf{+0.107}$\\
 & 0.25 & 1.345 & 1.333 & $+0.9$ & 0.014 & 0.023 & $+0.009$\\
 & 0.50 & 1.356 & 1.355 & $+0.1$ & 0.026 & 0.027 & $+0.001$\\
 & 0.75 & 0.942 & 0.940 & $+0.2$ & 0.027 & 0.029 & $+0.002$\\
 & 0.90 & 0.491 & 0.483 & $+1.7$ & 0.085 & 0.100 & $+0.015$\\
 & 0.95 & 0.291 & 0.274 & $+5.9$ & 0.161 & 0.210 & $+0.050$\\
\bottomrule
\end{tabular}
\begin{tablenotes}[flushleft]\footnotesize
\item \textit{Notes:} Highlighted rows = headline $(h{=}12,\ \tau{\leq}0.10)$ results.
Enhanced weakly dominates at every in-sample cell; minimum improvement $+0.10\%$.
Improvements concentrated in tails; median ($\tau{=}0.50$) improvements near zero at all horizons.
In-sample, 1997M01--2025M12 (effective sample for $h{=}12$: 1997M01--2024M12).
\end{tablenotes}
\end{threeparttable}
\end{center}\end{table}

\FloatBarrier

% ═══════════════════════════════════════════════════════════════════════════════
\section{Layer 0: Real-Time Event Detection}\label{sec:layer0}
% ═══════════════════════════════════════════════════════════════════════════════

\subsection{GDELT 2.0 News Surveillance}

The Global Database of Events, Language, and Tone (GDELT 2.0,
\citealp{leetaru2013}) monitors broadcast, print, and online news
in over 100 languages, updated every 15 minutes.
We query the GDELT Doc API~v2 with a Boolean AND of Middle East location
terms (Iran, Israel, Gaza, Lebanon, Yemen, Strait of Hormuz, Red Sea,
Persian Gulf, Saudi Arabia) crossed with conflict vocabulary
(military, strike, attack, missile, bomb, naval, escalation, drone).
The lookback window is 48 hours in calm regime, extended to 168 hours
in elevated or crisis regime.
Up to 250 articles are returned, sorted by recency.
Each article is scored for location relevance and CAMEO conflict severity;
the top-20 by combined score are passed to the LLM classifier.

\subsection{LLM CAMEO Coding}

The 20 ranked headlines are sent to an LLM (Anthropic Claude) with the
CAMEO severity guide \citep{schrodt2009}, the current GPR $z$-score as
an anchor, and instructions to return validated JSON with fields:
\texttt{severity}~(0--10), \texttt{cameo\_codes}~(list),
\texttt{gpr\_z\_estimate}, \texttt{deescalation}~(bool),
and \texttt{summary} ($\leq 100$ words).
A regex fallback operates on API unavailability.
The composite severity score is:
\begin{equation}
S_t = 0.5\,s^{\text{LLM}}_t + 0.3\,s^{\text{rule}}_{t,P_{75}}
    + 0.2\,\mathbf{1}[\text{AIS anomaly}_t].
\label{eq:severity}
\end{equation}
Weights sum to unity by construction.
The AIS component $\mathbf{1}[\text{AIS anomaly}_t]$ activates when
the freight, tanker equity, or headline-override conditions described
in Section~\ref{sec:ais} are met.
The $P_{75}$ of rule-based scores (rather than mean or max) is used
as a robust measure of event severity that discounts spurious
low-severity article matches.
Sensitivity: varying weights across
$(0.6,\,0.2,\,0.2)$, $(0.4,\,0.4,\,0.2)$, and $(0.33,\,0.33,\,0.33)$
shifts the March~2026 nowcast $\hat{z}_t$ by at most $\pm 0.12\sigma$,
leaving the \textsc{elevated} classification unchanged.

\subsection{AIS Tanker and Freight Proxy}\label{sec:ais}

We proxy Strait of Hormuz disruption using three complementary signals:
(i)~four listed VLCC operators (INSW, TK, TRMD, FRO),
(ii)~the Breakwave Tanker Shipping ETF (BWET), which holds 90\%
TD3C VLCC futures on the Middle East Gulf--to--Asia route and
provides the most direct market proxy for Hormuz freight rates,
and (iii)~the Brent--WTI spread.
All $z$-scores use a 90-day rolling window to avoid normalising
sustained disruptions.
The anomaly flag fires under any of three conditions:
(a)~BWET freight $z > 1.5$;
(b)~tanker basket $z > 1.5$ combined with Brent--WTI $z > 0.5$;
or (c)~both the tanker basket and Brent--WTI spread $z$-scores
simultaneously exceed $1.5\sigma$.
A headline-override mechanism supplements the market signals:
when GDELT headlines confirm a Hormuz closure or blockade and
CAMEO codes in the 19x range (impose restrictions) are active,
the anomaly flag is forced on regardless of market data,
guarding against situations where financial data lags physical
disruptions.

\subsection{Bridging the Publication Lag}
\label{sec:lag}

The GPR index is published with a one-month lag: as of March 2026, the
most recent official release covers November 2025.
This creates a blind spot for real-time risk monitoring---precisely the
period in which fast-moving geopolitical developments are most
policy-relevant.

Layer~0 bridges this gap by producing a \textit{sub-monthly GPR nowcast}
$\hat{z}_t$ derived from live GDELT headlines, LLM CAMEO severity scoring,
and the AIS tanker proxy.
Because GDELT ingests news continuously and the pipeline can be re-run
at any time, the signal is in principle updated daily---providing
actionable, timely risk signals precisely when they are most needed,
rather than waiting for the next official monthly release.
The nowcast is injected into the GaR model as the current-month
$\text{GPR}_z$ observation, replacing the stale monthly reading with a
contemporaneous estimate.
Formally, define the publication lag as $\ell = 1$ month.
The nowcast augments the information set from
$\mathcal{I}_{t-\ell}$ (monthly GPR only) to
$\mathcal{I}_t$ (monthly GPR + real-time Layer~0 signal),
so the enhanced GaR forecast conditions on
\begin{equation}
\hat{z}_t \;=\; 0.6\,S_t \;+\; 0.4\,\hat{z}^{\text{GPR}}_{t-\ell},
\label{eq:nowcast}
\end{equation}
where $S_t$ is rescaled from its 0--10 range to GPR $z$-score units
via $S_t / (10/4) = S_t \cdot 0.4$, so that a maximum-severity event
maps to approximately $z = 4$.\footnote{%
  For the mid-March reading: $S_t = 5.52$ maps to $5.52 \times 0.4 = 2.21$
  in $z$-scale; the nowcast is $0.6 \times 2.21 + 0.4 \times 2.65 = 2.39$.
  The reported $\hat{z} = 1.97$ is lower because the mid-March run used
  the regex rule-based CAMEO fallback (LLM unavailable), which assigns
  lower per-article severity scores than the LLM classifier, reducing $S_t$
  in the weighted average relative to the $z$-anchor.}
$S_t$ already incorporates the AIS anomaly flag with weight 0.2
(Equation~\ref{eq:severity}), so including AIS separately would
double-count its contribution; $\hat{z}^{\text{GPR}}_{t-\ell}$ is the
last official GPR $z$-score providing a mean-reverting anchor.

The Layer~0 pipeline provides two snapshots that bracket the March~2026
Hormuz crisis, illustrating the framework's capacity to detect a regime
transition in real time.

\paragraph{Mid-March (16~March 2026).}
The pipeline returned an \textsc{elevated} signal ($\hat{z} = 1.97$,
composite severity $S_t = 5.52/10$) under a rule-based CAMEO fallback.
Active CAMEO codes were 191 (``Impose blockade, restrict movement'')
and 15 (``Exhibit military posture or force buildup''),
reflecting regional force movements and access restrictions.
Despite this elevated backdrop, the AIS trigger remained inactive:
the tanker equity basket traded at $z = -0.68$, BWET freight at $z = -0.42$,
and the Brent--WTI spread stood at \$0.70 ($z = -3.38$)---
the oil market was assigning near-zero probability of a Hormuz
supply disruption.
GIPI stood at $-1.01$ (December 2025), confirming that all five
inflation transmission channels remained subdued.
This was the paper's central thesis in real time: military posture
activity was elevated, yet the energy and supply-chain channels through
which GPR transmits to inflation and output were closed.
The six-month GaR enhanced model nowcast showed a median forecast of
$+1.30\%$, recession probability 2.7\%, and conditional skewness $+0.27$:
a benign central case with upside skew, buffered by subdued GIPI.

\paragraph{Late March (25~March 2026).}
Following Iran's closure of the Strait of Hormuz---which reduced transit
traffic by approximately 70\% and pushed Brent crude above \$126---the
pipeline escalated to a \textsc{crisis} signal ($\hat{z} = 2.80$,
composite severity $S_t = 4.8/10$, LLM-coded).
Active CAMEO codes shifted to 195 (``Employ aerial weapons''),
204 (``Occupy territory''), and 191 (``Impose blockade''),
reflecting the Israel--Iran military escalation and reciprocal strikes.
The headline-override mechanism activated: GDELT headlines confirming
the Hormuz closure, combined with CAMEO 19x codes, forced the AIS
anomaly flag on despite lagging market data (tanker basket $z = -1.54$,
BWET freight $z = -1.30$, Brent--WTI spread \$7.40 with $z = +0.84$).
This late-March reading demonstrates precisely the scenario the paper's
framework is designed to detect: the transition from ``fire without
smoke'' (elevated GPR, closed channels) to active supply-channel
transmission as a genuine chokepoint disruption begins to reopen the
inflation channels.

\subsection{Current Nowcast}

Table~\ref{tab:l0} presents both the mid-March and late-March signals
side by side, documenting the regime transition numerically.
The full interactive Layer~0 panel is available at
\texttt{https://geoshock.onrender.com}.

% Note: The Layer~0 panel is available interactively at
% https://geoshock.onrender.com; the numerical content is
% reported in Table~\ref{tab:l0}.

\begin{table}[htbp]
\caption{Layer~0 Real-Time Signal: Mid-March vs.\ Late March 2026}
\label{tab:l0}
\begin{center}
\begin{threeparttable}\small\setlength{\tabcolsep}{5pt}
\begin{tabular}{@{}p{4.0cm} p{3.0cm} p{3.0cm} p{3.5cm}@{}}
\toprule
\textit{Metric} & \textit{16 Mar} & \textit{25 Mar} & \textit{Change} \\
\midrule
\multicolumn{4}{l}{\textit{A. Real-time event signal}}\\[2pt]
Regime & \textsc{elevated} & \textsc{crisis} & Escalation\\
Composite severity $S_t$ & 5.52\,/\,10 & 4.8\,/\,10 & ---$^{\dagger}$\\
GPR nowcast $\hat{z}$ & 1.97 & 2.80 & $+0.83$\\
CAMEO codes & 191, 15 & 195, 204, 191 & Kinetic codes active\\
GDELT articles & 3 (100\% ME) & 50 (44\% ME) & Volume surge\\
AIS anomaly flag & \textbf{No} & \textbf{Yes}$^{\ddagger}$ & Headline override\\
Tanker basket $z$ & $-0.68$ & $-1.54$ & ---\\
BWET freight $z$ & $-0.42$ & $-1.30$ & ---\\
Brent--WTI spread & \$0.70 & \$7.40 & Widening\\
\midrule
\multicolumn{4}{l}{\textit{B. Monthly anchors (most recent releases)}}\\[2pt]
GPR level (Nov 2025) & \multicolumn{2}{c}{371.1, $z=2.65$} & \textsc{crisis} band\\
GIPI (Dec 2025) & \multicolumn{2}{c}{$-1.01$} & Subdued\\
\midrule
\multicolumn{4}{l}{\textit{C. Implied 6-month GaR nowcast (enhanced model)}}\\[2pt]
Median IP YoY & $+1.3\%$ & \multicolumn{2}{l}{To be updated as GIPI revises}\\
GaR$_5$ & $-3.2\%$ & \multicolumn{2}{l}{Left tail expected to widen}\\
$\Pr(\text{recession})$ & 2.7\% & \multicolumn{2}{l}{Rising if channels reopen}\\
\bottomrule
\end{tabular}
\begin{tablenotes}[flushleft]\footnotesize
\item $^{\dagger}$Composite severity $S_t$ fell slightly (5.52 $\to$ 4.8)
because the 25~Mar article pool was broader (50 articles, 44\% ME-relevant
vs.\ 3 articles, 100\% ME), diluting the $P_{75}$ rule-based severity
component (Equation~\ref{eq:severity}).
The GPR nowcast $\hat{z}$ nevertheless rose sharply (1.97 $\to$ 2.80)
because the LLM severity estimate and the GPR $z$-anchor both increased.
This divergence illustrates a feature of the composite design:
$S_t$ reflects per-article severity concentration, while $\hat{z}$
captures the overall geopolitical state.
\item $^{\ddagger}$Headline override: GDELT headlines confirmed Hormuz closure
and CAMEO 19x codes (blockade) were active, forcing the AIS anomaly flag
despite lagging market data.
\item \textit{Notes:}
CAMEO code 15 = military posture; 191 = impose blockade;
195 = employ aerial weapons; 204 = occupy territory.
The shift from codes 191/15 to 195/204/191 marks the transition from
posturing to kinetic operations.
AIS proxy uses 90-day rolling $z$-scores on VLCC equity returns,
BWET freight (TD3C route), and Brent--WTI spread.
The mid-March reading illustrates the paper's central thesis
(fire without smoke); the late-March reading illustrates the regime
transition as the Strait closure begins to reopen transmission channels.
\end{tablenotes}
\end{threeparttable}
\end{center}
\end{table}

The two snapshots in Table~\ref{tab:l0} bracket a regime transition.
On 16~March, CAMEO codes 191 and 15 indicated access restrictions and
military posture, yet the AIS trigger was inactive and the Brent--WTI
spread stood at \$0.70---the oil market was not pricing a Hormuz disruption.
This was the paper's central thesis in real time: elevated geopolitical
activity with subdued GIPI ($-1.01$) produced an elevated GPR nowcast
but no supply-channel transmission.
By 25~March, the situation had changed materially: Iran's closure of the
Strait reduced transit traffic by approximately 70\% and pushed Brent
crude above \$126.
The pipeline detected this shift through two channels.
First, the GDELT article count surged from 3 to 50, with CAMEO codes
escalating from posturing (191, 15) to kinetic operations (195, 204).
Second, the headline-override mechanism activated: news confirming the
Hormuz closure, combined with CAMEO 19x blockade codes, forced the AIS
anomaly flag despite lagging market data.
The Brent--WTI spread widened from \$0.70 to \$7.40, consistent with
the early stages of a Hormuz premium emerging---analogous to (though
smaller than) the \$8+ spread observed after the September~2019 Abqaiq
attack.
This transition from ``fire without smoke'' to active supply-channel
disruption is precisely the scenario the GIPI framework is designed to
capture: as transmission channels reopen, the model's left-tail mass
is expected to expand, shifting the growth distribution toward higher
recession probability.

% ═══════════════════════════════════════════════════════════════════════════════
\section{Conclusion}\label{sec:conclude}
% ═══════════════════════════════════════════════════════════════════════════════

\noindent\textbf{What this paper does.}
We introduce GIPI, a monthly index of geopolitical inflation transmission
pressure---the first principal component of five channel variables---and
integrate it into a Growth-at-Risk framework via a GPR$\times$GIPI
interaction term.
The core idea is simple: a geopolitical shock is most dangerous to the
growth distribution when the inflation channels are open (high GIPI), and
least dangerous when they are closed (low GIPI, as in 2025).

\noindent\textbf{What we find.}
Four results emerge.
\textit{First}, the amplification mechanism is statistically robust
\textit{in sample}: GIPI carries independent left-tail information in every
16 of 18 quantile-horizon orthogonalisation cells in
Table~\ref{tab:orth}---including all left-tail cells ($\tau \leq 0.25$)---with
$\tau = 0.90$ scoring-rule improvements also positive (Table~\ref{tab:pinball}).
The orthogonalised coefficient at $h=12$, $\tau=0.05$ is $-2.10$
($p<0.01$)---the single strongest result in the paper.
\textit{Second}, conditioning on GIPI changes the in-sample shape of
the IP growth distribution: the 12-month, 5th-quantile cell achieves
14.7\% better in-sample fit and pseudo-$R^2$ rises from 0.146 to 0.271,
with LASSO identifying import price growth as the dominant long-horizon
channel.
\textit{Third}, the 2023--25 GPR--GIPI decoupling ($\rho=0.131$) is
unprecedented in the post-2000 sample: US shale production, IEA reserve
coordination, and Fed credibility have simultaneously closed all five
transmission channels.
The model's in-sample conditional distribution assigns a 2.7\% six-month
recession probability at current GIPI ($-1.01$), with positive skewness
($+0.27$) reflecting subdued channel pressure.
\textit{Fourth}, Layer~0 provides a real-time stress test across
March~2026.
On 16~March, the system illustrates the paper's thesis: CAMEO codes 191
and 15 are active with the GPR nowcast at $\hat{z}=1.97$, yet the AIS
trigger is inactive and Brent--WTI stands at \$0.70---``fire without smoke.''
By 25~March, following Iran's closure of the Strait of Hormuz (70\%
traffic reduction, Brent above \$126), the pipeline escalates to crisis
($\hat{z}=2.80$, CAMEO 195/204/191, headline-override AIS anomaly),
demonstrating the framework's capacity to detect regime transitions
as transmission channels reopen.

\noindent\textbf{When does the decoupling break?}
The mid-March in-sample distribution showed positive skewness ($+0.27$ at $h=6$):
upside scenarios dominated under subdued GIPI, yet GaR$_5 = -3.2\%$ documented
that a left tail remained if channels re-engaged.
The decoupling breaks down if \textit{any} of four conditions changes:
(i)~US production falls below 12\,mb/d (possibly from tariff-driven capex
cuts); (ii)~IEA reserve buffers are depleted; (iii)~inflation expectations
de-anchor, as measured by a GIPI reading above $+2$; or
(iv)~tariff-driven import-price pass-through reopens the cost-push
channel independently of energy---a scenario increasingly relevant
given 2025--26 US trade-policy volatility, which the current GIPI
captures partially through the import-price component but does not
separately identify from geopolitical supply shocks.
The late-March Hormuz closure suggests condition~(ii) may now be under stress:
a sustained Strait blockade would draw down strategic reserves and push
the GPR--GIPI correlation back toward its 2022 value, expanding the model's
left-tail mass sharply.
The Layer~0 pipeline provides 48-hour early warning via the AIS anomaly
trigger (now incorporating VLCC freight rates and headline override)
and the breakeven component of GIPI.

\noindent\textbf{Extensions.}
Four frontier extensions would deepen the framework.
\textit{First}, a real-time GIPI construction---re-estimating PCA
weights at each rolling origin and applying the Giacomini and White
\citeyear{giacomini2006} fixed-window test, which is valid for nested
model comparison under a rolling window design---would enable a fully
rigorous out-of-sample forecasting evaluation.
This is the natural next step for the research agenda opened by this paper.
\textit{Second}, formal instrument-based identification of GIPI---through
the exogenous oil supply news series of \citet{baumeister2019} or OPEC
spare-capacity shocks \citep{caldara2022}---would convert the current
conservative lower-bound estimates into point-identified structural
amplification effects.
The three endogeneity controls already employed (lagged GIPI, forward
projection design, orthogonalisation) are strong methodological
foundations for this extension.
\textit{Third}, real vessel-level AIS data (Kpler or Spire Maritime)
would supplement the current BWET freight and equity proxy for a
sharper Hormuz signal; the March~2026 headline-override mechanism
demonstrates the value of multi-source confirmation.
\textit{Fourth}, state-dependent local projections---estimated separately
in high-GIPI and low-GIPI regimes---would formally test whether the
amplification hypothesis operates symmetrically at the IRF level.
\textit{Fifth}, sensitivity analysis for GIPI construction---testing
an equally-weighted channel average against PCA, or retaining PC2
(eigenvalue $= 1.005$, marginally above the Kaiser criterion) alongside
PC1---would strengthen the robustness of the index design.
\textit{Sixth}, distinguishing the tariff-driven import-price channel
from geopolitical supply shocks within GIPI, and incorporating
the shortage index of \citet{caldara2025shortages} as a complementary
channel-openness measure, would sharpen the framework's ability
to separate concurrent transmission mechanisms.

\newpage

% ── References ────────────────────────────────────────────────────────────────
\bibliographystyle{aer}
\begin{thebibliography}{99}

\bibitem[Bernanke and Blanchard(2023)]{bernanke2023}
Bernanke, B.S., and O.~Blanchard. 2023.
``What Caused the US Pandemic-Era Inflation?''
\textit{Brookings Papers on Economic Activity}, Spring 2023: 1--86.

\bibitem[Adrian et al.(2019)]{adrian2019}
Adrian, T., N.~Boyarchenko, and D.~Giannone. 2019.
``Vulnerable Growth.'' \textit{American Economic Review}
109(4): 1263--1289.

\bibitem[Benigno et al.(2022)]{benigno2022}
Benigno, G., J.~di Giovanni, J.J.J.~Groen, and A.I.~Noble. 2022.
``The GSCPI: A New Barometer of Global Supply Chain Pressures.''
FRBNY Staff Reports No.\ 1017.

\bibitem[Bloom(2009)]{bloom2009}
Bloom, N. 2009. ``The Impact of Uncertainty Shocks.''
\textit{Econometrica} 77(3): 623--685.

\bibitem[Bonadio et al.(2021)]{bonadio2021}
Bonadio, B., Z.~Huo, A.A.~Levchenko, and N.~Pandalai-Nayar. 2021.
``Global Supply Chains in the Pandemic.''
\textit{Journal of International Economics} 133: 103534.

\bibitem[Caldara and Iacoviello(2022)]{caldara2022}
Caldara, D., and M.~Iacoviello. 2022.
``Measuring Geopolitical Risk.''
\textit{American Economic Review} 112(4): 1194--1225.

\bibitem[Caldara et al.(2025)]{caldara2025}
Caldara, D., S.~Conlisk, M.~Iacoviello, and M.~Penn. 2025.
``Do Geopolitical Risks Raise or Lower Inflation?''
\textit{Journal of Monetary Economics}, forthcoming.

\bibitem[Carri\`{e}re-Swallow et al.(2023)]{carriereswallow2023}
Carri\`{e}re-Swallow, Y., N.C.~Firat, D.~Furceri, and D.~Jimenez. 2023.
``State-Dependent Exchange Rate Pass-Through.''
\textit{Journal of International Money and Finance} 134: 102830.

\bibitem[Jord\`{a}(2005)]{jorda2005}
Jord\`{a}, \`{O}. 2005.
``Estimation and Inference of Impulse Responses by Local Projections.''
\textit{American Economic Review} 95(1): 161--182.

\bibitem[Kilian(2009)]{kilian2009}
Kilian, L. 2009.
``Not All Oil Price Shocks Are Alike.''
\textit{American Economic Review} 99(3): 1053--1069.

\bibitem[Koenker and Bassett(1978)]{koenker1978}
Koenker, R., and G.~Bassett. 1978.
``Regression Quantiles.'' \textit{Econometrica} 46(1): 33--50.

\bibitem[Koenker and Machado(1999)]{koenker1999}
Koenker, R., and J.A.F.~Machado. 1999.
``Goodness of Fit and Related Inference Processes for Quantile Regression.''
\textit{Journal of the American Statistical Association}
94(448): 1296--1310.

\bibitem[Leetaru and Schrodt(2013)]{leetaru2013}
Leetaru, K., and P.A.~Schrodt. 2013.
``GDELT: Global Data on Events, Location, and Tone, 1979--2012.''
\textit{ISA Annual Conference}. San Francisco.

\bibitem[L\'{o}pez-Salido and Loria(2024)]{lopezloria2024}
L\'{o}pez-Salido, D., and F.~Loria. 2024.
``Inflation at Risk.''
\textit{Journal of Monetary Economics} 145: 103568.

\bibitem[Newey and West(1987)]{newey1987}
Newey, W.K., and K.D.~West. 1987.
``A Simple, Positive Semi-definite, Heteroskedasticity and
Autocorrelation Consistent Covariance Matrix.''
\textit{Econometrica} 55(3): 703--708.

\bibitem[Schrodt(2009)]{schrodt2009}
Schrodt, P.A. 2009.
``CAMEO: Conflict and Mediation Event Observations Codebook.''
Version 1.1b3. Pennsylvania State University.

\bibitem[Conflitti and Luciani(2019)]{conflitti2019}
Conflitti, C., and M.~Luciani. 2019.
``Oil Price Pass-Through into Core Inflation.''
\textit{The Energy Journal} 40(6): 221--248.


\bibitem[Adams et al.(2021)]{adams2021}
Adams, P.A., T.~Adrian, N.~Boyarchenko, and D.~Giannone. 2021.
``Forecasting Macroeconomic Risks.''
\textit{International Journal of Forecasting} 37(3): 1173--1191.

\bibitem[Aiyar et al.(2023)]{aiyar2023}
Aiyar, S., A.~Ilyina, J.~Chen, A.~Ebeke, R.~Garcia-Saltos,
T.~Gueorguiev, A.~Mangal, T.~Nhung Nguyen, D.J.~Oner,
F.~Raei, M.~Srinivasan, and H.~Tokuoka. 2023.
``Geoeconomic Fragmentation and the Future of Multilateralism.''
IMF Staff Discussion Note SDN/2023/001.

\bibitem[Baker, Bloom and Davis(2016)]{baker2016}
Baker, S.R., N.~Bloom, and S.J.~Davis. 2016.
``Measuring Economic Policy Uncertainty.''
\textit{Quarterly Journal of Economics} 131(4): 1593--1636.

\bibitem[Baumeister and Hamilton(2019)]{baumeister2019}
Baumeister, C., and J.D.~Hamilton. 2019.
``Structural Interpretation of Vector Autoregressions with Incomplete
Identification: Revisiting the Role of Oil Supply and Demand Shocks.''
\textit{American Economic Review} 109(5): 1873--1910.

\bibitem[Caldara, Iacoviello and Yu(2025)]{caldara2025shortages}
Caldara, D., M.~Iacoviello, and D.~Yu. 2025.
``Measuring Shortages since 1900.''
International Finance Discussion Papers No.~1407,
Board of Governors of the Federal Reserve System.

\bibitem[Chavleishvili and Manganelli(2024)]{chavleishvili2024}
Chavleishvili, S., and S.~Manganelli. 2024.
``Forecasting and Stress Testing with Quantile Vector Autoregression.''
\textit{Journal of Applied Econometrics} 39(1): 66--85.

\bibitem[Federle et al.(2024)]{federle2024}
Federle, J., A.~Meier, G.J.~M\"{u}ller, W.~Mutschler, and M.~Schularick. 2024.
``The Price of War.''
CEPR Discussion Paper No.~18834 / Kiel Working Paper No.~2262.

\bibitem[Fern\'{a}ndez-Villaverde, Mineyama and Song(2024)]{fernandez2024}
Fern\'{a}ndez-Villaverde, J., T.~Mineyama, and D.~Song. 2024.
``Are We Fragmented Yet? Measuring Geopolitical Fragmentation and Its
Causal Effects.'' NBER Working Paper No.~32638.

\bibitem[Hassan et al.(2019)]{hassan2019}
Hassan, T.A., S.~Hollander, L.~van~Lent, and A.~Tahoun. 2019.
``Firm-Level Political Risk: Measurement and Effects.''
\textit{Quarterly Journal of Economics} 134(4): 2135--2202.

\bibitem[Metiu(2025)]{metiu2025}
Metiu, N. 2025.
``Global Financial Transmission of Geopolitical Risk Shocks.''
Deutsche Bundesbank / SSRN Working Paper No.~5195860.

\bibitem[Pinchetti(2024)]{pinchetti2024}
Pinchetti, M. 2024.
``Geopolitical Risk and Inflation: The Role of Energy Markets.''
Centre for Macroeconomics Discussion Paper No.~2431.
(Also Banque de France Working Paper No.~1005, 2025.)

\bibitem[Plagborg-M{\o}ller and Wolf(2021)]{plagborg2021}
Plagborg-M{\o}ller, M., and C.K.~Wolf. 2021.
``Local Projections and VARs Estimate the Same Impulse Responses.''
\textit{Econometrica} 89(2): 955--980.




\bibitem[{EIA}(2025)]{eia2025}
U.S. Energy Information Administration. 2025.
\textit{Monthly Energy Review}, Table~3.1: Petroleum Overview.
December 2025. Washington, DC: EIA.

\bibitem[Giacomini and White(2006)]{giacomini2006}
Giacomini, R., and H.~White. 2006.
``Tests of Conditional Predictive Ability.''
\textit{Econometrica} 74(6): 1545--1578.

\end{thebibliography}

\end{document}
